\section{Síntesis y ampliación}

\subsection{Una tabla que recoge toda la Recipe}

La siguiente tabla reúne los cuatro sistemas principales según su entrada, su salida y sus modos de falla. Es más útil que memorizarlos por año de publicación, porque cualquier robot foundation model posterior debe cubrir estas funciones, aunque su implementación pueda integrarse en un mismo modelo.

\begin{longtable}{@{}L{0.16\textwidth}L{0.22\textwidth}L{0.23\textwidth}L{0.27\textwidth}@{}}
\toprule
Sistema & Entrada principal & Salida principal & Modo de falla principal \\
\midrule
RT-1 & Seis cuadros visuales, tarea en lenguaje, demostraciones fuera de línea & Acciones discretizadas en tiempo real & Escenas nuevas, variación de varios factores a la vez, context de largo plazo y dataset coverage \\
SayCan & Objetivo de alto nivel, skill list, affordance scores & Siguiente grounded skill & Skill library demasiado pequeña, affordance imprecisa, fallas de ejecución \\
Inner Monologue & Scene, success, retroalimentación del usuario, planes previos & Belief actualizada y plan del siguiente paso & Errores del detector / caption, ontology de retroalimentación incompleta \\
DIAL & Poco lenguaje humano, muchas offline trajectories, VLM & Language labels y policy ampliadas & Label noise, cobertura insuficiente de comportamientos, evidencia limitada de generalización composicional \\
\bottomrule
\end{longtable}

\subsection{Conclusiones finales}

Primero, el recurso central del scaling en robótica no es un único episode count, sino la cobertura diversa de tareas, entornos, lenguaje y embodiment. Segundo, uno de los valores reales del Transformer es la token interface y un sequence model de alta capacidad, pero el presupuesto de tiempo real obliga a comprimir la visión, usar un historial corto y desplegar modelos pequeños. Tercero, la mejor forma de conectar un LLM a un robot es mediante un planner restringido por grounding, no tratando directamente las probabilidades del lenguaje como acciones. Cuarto, el feedback y el VLM relabeling pueden inyectar en el sistema el prior de foundation models externos, pero no pueden crear de la nada habilidades físicas que no se hayan recolectado.

\begin{importantbox}{La frase que más vale llevarse}
La robotics foundation-model recipe de 2023 no es «conectar el robot a un LLM», sino: entrenar con datos fuera de línea grandes y diversos un skill model ejecutable en tiempo real, usar el lenguaje para conectar la planeación con la retroalimentación, usar la affordance para mantener el grounding y dejar que foundation models externos mejoren de forma continua la interfaz semántica.
\end{importantbox}

\subsection{Lecturas adicionales}

Todos los materiales siguientes se publicaron antes de la fecha de la clase y forman directamente la cadena de evidencia de esta sesión. El orden de lectura sugerido es: la capa de skill/data de RT-1, el planner de SayCan, el feedback de Inner Monologue y, por último, DIAL para volver a la interfaz de datos.

\begin{itemize}
  \item \href{https://arxiv.org/abs/2212.06817}{Brohan et al., RT-1: Robotics Transformer for Real-World Control at Scale}: imitación multitarea, arquitectura, data scaling y ablation.
  \item \href{https://arxiv.org/abs/2204.01691}{Ahn et al., Do As I Can, Not As I Say}: LLM usefulness y affordance grounding.
  \item \href{https://arxiv.org/abs/2207.05608}{Huang et al., Inner Monologue}: planeación en lazo cerrado con scene, success y active feedback.
  \item \href{https://arxiv.org/abs/2202.02005}{Jang et al., BC-Z}: imitación multitarea condicionada por lenguaje y zero-shot task generalization.
  \item \href{https://arxiv.org/abs/2211.11736}{Xiao et al., Robotic Skill Acquisition via Instruction Augmentation with Vision-Language Models}: DIAL y VLM relabeling.
\end{itemize}

\subsection{Autoevaluación posterior a la clase}

\begin{importantbox}{Diez preguntas de autoevaluación}
\begin{enumerate}
  \item ¿Por qué esta sesión solo puede decir «towards a robotics foundation model» y no que ya se demostró la emergence?
  \item ¿Cómo desacopla el offline robot learning la data generation de la consumption, y qué riesgo de distribution-shift introduce?
  \item ¿Qué restricción de recursos resuelve cada uno de estos elementos de RT-1: TokenLearner, el historial de seis cuadros y los 256 bins?
  \item ¿Por qué reducir la task diversity perjudica más la generalization que una episode reduction equivalente?
  \item En SayCan, ¿qué es lo que el LM score y el affordance score no pueden sustituir, cada uno?
  \item ¿Por qué deben reportarse por separado el planning success y el execution success?
  \item ¿En qué se diferencian el passive, el active y el success feedback de Inner Monologue?
  \item ¿Por qué DIAL puede ampliar la cobertura del lenguaje, pero no crear habilidades que no se recolectaron?
  \item ¿De qué condiciones del dataset depende «poco label noise: okay», y de qué manera no debe malinterpretarse?
  \item Si seis cuadros ya se acercan al context limit, ¿qué memory architecture necesitan las tareas que duran minutos?
\end{enumerate}
\end{importantbox}

\end{document}
